\section{Endpoint ledger and the next coefficient-specific gate}
\label{sec:tpc35-ledger-next-gate}

The results above have different logical status.  The Gram lower bounds
hold for abstract phase classes, the projective saturation theorem concerns
a particular right--hand side, the character formula is an exact finite
diagonalization, the CRT ladder is exact arithmetic bookkeeping, and the
same--ultra estimate is a physical upper bound.  We compare them without
turning a method limitation into a physical lower bound.

\subsection{Energy ledger}

At
\[
 Q=X^{267/400+o(1)},\qquad
 J=X^{133/400+o(1)},\qquad
 T=X^{193/500+o(1)},
\]
the orbit-energy scales are as follows.

\begin{table}[H]
\centering
\caption{Orbit-sliced energy ledger at the high--\(\beta\) endpoint.}
\label{tab:tpc35-energy-ledger}
\small
\begin{tabular}{@{}>{\raggedright\arraybackslash}p{0.39\textwidth}
>{\raggedright\arraybackslash}p{0.25\textwidth}
>{\raggedleft\arraybackslash}p{0.20\textwidth}@{}}
\toprule
Object & Scale & \(X\)-exponent \\
\midrule
Absolute orbit-sliced baseline & \(Q^3J\) & \(4670/2000\) \\
One-modulus multiplicity certificate, optimistically applied
 & \(Q^3\) & \(4005/2000\) \\
Physical same-ultra block & \(Q^3J/T\) & \(3898/2000\) \\
Sufficient orbit-sliced gate & \(Q^3/J\) & \(3340/2000\) \\
Proved orbit-sliced diagonal & \(Q^2J\) & \(3335/2000\) \\
Three-modulus aliases times diagonal & \((X/J^3)Q^2J\) & \(3340/2000\) \\
\bottomrule
\end{tabular}
\end{table}

The last row is conditional accounting: it assumes a diagonal-scale bound
for every alias fiber.  The equality of its exponent with the gate is the
critical match \eqref{eq:tpc35-critical-three-modulus}, not a proved
estimate.

The uniform norm comparison is even more decisive within the abstract
phase class.

\begin{table}[H]
\centering
\caption{Uniform Gram scales in the equal-deletion phase model.}
\label{tab:tpc35-uniform-ledger}
\small
\begin{tabular}{@{}llll@{}}
\toprule
Quantity & Abstract lower & Required scale & Ratio \\
\midrule
Operator norm & \(Q^2\) & \(Q^2/J\) & \(J\) \\
Hilbert--Schmidt norm squared & \(Q^4J\) & \(Q^4/J^2\) & \(J^3\) \\
\bottomrule
\end{tabular}
\end{table}

These ratios explain why replacing the physical coefficient direction by
all unit vectors is too expensive in a broad model.  They do not determine
the norms of the actual masked matrices.

\subsection{Two viable next formulations}

The exact product spectrum suggests a signed finite-field target.  On a
layer already reduced to the complete separable form
\eqref{eq:tpc35-product-operator}, the quantity to control is
\begin{equation}
 \mathfrak E_q(b;f,g)
 :=\frac1{q-1}\sum_\chi
 |\widehat b(\bar\chi)|^2\cE_\chi,
 \label{eq:tpc35-signed-spectral-target}
\end{equation}
By \cref{thm:tpc35-product-spectrum}, this is exactly the coherent energy
of that finite layer.  To feed it into
\eqref{eq:tpc35-coefficient-gate}, one must first decompose the actual
\(j\)-dependent mask and incomplete ranges into such layers with total
cost \(X^{o(1)}\).  That physical interface lemma is not proved here.
After it is supplied, a theorem showing that the resulting
\(\widehat b\)'s avoid characters with large \(\cE_\chi\) would be
sufficient; it need not bound \(\max_\chi\cE_\chi\) for every abstract
coefficient vector.

The CRT formulation keeps three pairwise coprime primes
\(q_1,q_2,q_3\asymp J\) separate.  Its character expansion has eight
coordinate faces:
\begin{equation}
 \varnothing,\quad
 \{1\},\{2\},\{3\},\quad
 \{1,2\},\{1,3\},\{2,3\},\quad
 \{1,2,3\}.
 \label{eq:tpc35-three-coordinate-faces}
\end{equation}
The empty face is the principal term; singleton and double faces are the
non--tensor modes exposed by \cref{thm:tpc35-CRT-lift}; the full face is
where three independent character oscillations may multiply.  A genuine
joint theorem must separately control or eliminate the principal face;
ordinary additive centering does not do so automatically, as
\eqref{eq:tpc35-principal-energy} shows.  The proper faces must likewise
be small, absorbed into previously bounded diagonals, or cancelled by the
actual coefficients.  Estimating only the full face is insufficient.

These are two descriptions of the same strategic requirement:
do not replace the actual coefficient direction by a maximum over all
characters, and do not replace joint CRT sampling by three independent
applications of a one--modulus mean square.

\subsection{Route menu}

The next step may proceed along any of the following nonexclusive routes.
\begin{enumerate}[label=\textup{Route \Alph*:},leftmargin=5.8em]
 \item \emph{Coefficient-weighted multiplicative spectrum.}
 First prove a controlled complete separable decomposition of the physical
 mask, then establish a dyadically averaged bound for
 \eqref{eq:tpc35-signed-spectral-target}, retaining the M\"obius row
 coefficients, prime-source restriction, and actual gcd projector.  It is
 enough to control spectral overlap, not every character separately.

 \item \emph{Joint three-modulus dispersion.}
 Expand the three CRT coordinates simultaneously, center before
 separating them, and prove diagonal-scale reassembly for all aliases and
 proper-coordinate modes.  The exact available alias budget is
 \(X/J^3=Q/J^2\); four moduli may be used for no--wrap, but they add a
 further coordinate face that also requires control.

 \item \emph{Terminal/proper signed coupling.}
 Use the terminal equality result to remove the same-divisor terminal
 off--diagonal, disperse the proper-divisor layers through their
 complementary factors, and retain cancellation between terminal and
 proper layers.  A triangle inequality separating them may destroy the
 needed cancellation.

 \item \emph{Hybrid incidence--spectrum.}
 Use short-orbit uniqueness and gcd/lcm pruning to lower the multiplicity
 of the bad character fibers, then apply the coefficient-specific
 spectrum only to the remaining high-degree points.  The ambient Schur
 bound \eqref{eq:tpc35-physical-incidence-scale} shows that incidence
 without signed spectral input is insufficient.
\end{enumerate}

The first two routes are now precise enough that failure can also be made
mathematical.  A construction concentrating the physical coefficient
mass on one bad character, or a lower-coordinate CRT mode surviving all
available centerings, would be a scoped negative result.  Conversely, a
positive estimate for only a restricted range of projector divisors or
alias fibers would still close a nontrivial subpacket.

\subsection{Relation to the preceding chain}

TPC-19 supplies primitive determinant--divisor CRT normal forms
\citep{WangTPC19}.  TPC-29 already proves the general reflected
compatibility \(g\mid h(m_1-m_2)\), its primitive specialization, and the
unique lcm class \citep{WangTPC29}.  The self-contained restatement in
\cref{prop:tpc35-affine-CRT-compatibility} adds the reduced criterion
\(g/(g,h)\mid m_1-m_2\) and records a counterexample to a tempting swapped
coprimality hypothesis; the compatibility theorem itself is not claimed
as new.  TPC-22 and TPC-24 already show that low conductor
coordinates can survive a product decomposition
\citep{WangTPC22,WangTPC24}; the new use here is at orbit-scale projective
moduli with the explicit three/four threshold.  TPC-33 proves the
unweighted complete product--slope lift and fixed-determinant projective
Plancherel identity \citep{WangTPC33}.  The new spectral step is
\cref{thm:tpc35-product-spectrum}, which diagonalizes a coherent weighted
product--slope sum.  TPC-34 first states the coefficient-specific orbit
gate, proves its diagonal, and gives a coherent constant-phase witness at
the sign-blind scale \citep{WangTPC34}.  The present rank--trace theorem
upgrades that witness to a lower bound for every unit-phase
equal-deletion frame, with Fourier sharpness.  The new physical step is
\cref{thm:tpc35-shared-ultra}, which proves a \(T\)-saving bound for the
opened same-ultra block after orbit slicing.

Even a future proof of \eqref{eq:tpc35-coefficient-gate} would close only
the inherited matched residual shell.  Earlier packet reconstruction,
main-term identification, and positivity stages would remain.  The word
``dynamical'' here refers to deterministic affine traversal
\(j\mapsto mj+h\) modulo divisor moduli.  No mixing theorem, logistic-map
conjugacy, M\"obius disjointness theorem, or spectral interpretation of
zeta zeros is asserted.
